\section{Despliegue}

La plataforma web se encuentra publicada en Vercel y utiliza Supabase
para la persistencia de información. La interfaz desarrollada con Angular
y la API implementada con Node.js y Express se presentan bajo un mismo
dominio, permitiendo acceder al sistema mediante un navegador web.

\subsection{Plataforma de alojamiento}

La configuración del proyecto se mantiene en el archivo
\path{vercel.json}. Este archivo define un servicio para el frontend,
cuya raíz es \path{frontend/}, y otro para el backend, ubicado en
\path{backend/}, con entrada en \path{src/app.js}.

Las solicitudes cuyo recorrido comienza con \path{/api} se dirigen
al servicio del backend; las demás se dirigen al frontend. Esta
distribución permite mantener la interfaz y los servicios HTTP bajo
la misma dirección pública. PostgreSQL en Supabase conserva los registros
estructurados y la integración con Supabase Storage proporciona el
almacenamiento de archivos radiográficos.

\subsection{Procedimiento de despliegue}

El procedimiento para publicar una revisión del sistema comprende
las siguientes actividades:

\begin{enumerate}
    \item Preparar la revisión del código fuente y conservarla en el
    repositorio Git del proyecto.

    \item Instalar las dependencias definidas en los espacios de trabajo
    del frontend y del backend. Comprobar la compilación de Angular
    mediante \texttt{npm run build} y ejecutar las pruebas correspondientes.

    \item Preparar la base de datos mediante los scripts SQL y las
    migraciones del proyecto, respetando el historial de cambios
    y evitando repetir la instalación inicial sobre una base existente.

    \item Configurar en Vercel las variables de entorno necesarias
    para la conexión con PostgreSQL, la firma de sesiones y el acceso
    al almacenamiento. Mantener el bucket de radiografías privado.

    \item Desplegar la revisión utilizando la configuración de servicios
    y rutas de \path{vercel.json}.

    \item Verificar la dirección pública, la salud de la API, el inicio
    de sesión y las operaciones protegidas. Para cerrar la validación
    del flujo radiográfico, comprobar también la solicitud, el enlace
    temporal, la entrega del archivo y su consulta autorizada.
\end{enumerate}

\subsection{Configuración de entornos}

Durante el desarrollo, Angular y Express se ejecutan en procesos locales
separados. El backend obtiene su configuración del archivo
\path{backend/.env}, excluido del repositorio mediante
\path{.gitignore}. La aplicación Angular utiliza un proxy para comunicar
las solicitudes de la interfaz con la API local.

La configuración de producción utiliza las variables de entorno del
servicio desplegado. Los parámetros relevantes incluyen
\texttt{NODE\_ENV}, \texttt{APP\_URL}, \texttt{DATABASE\_URL},
\texttt{JWT\_SECRET}, \texttt{SUPABASE\_URL},
\texttt{SUPABASE\_SECRET\_KEY} y
\texttt{SUPABASE\_STORAGE\_BUCKET}.

Para producción se establece \texttt{NODE\_ENV=production},
\texttt{LOCAL\_ACCESS=false} y \texttt{ENFORCE\_ROLES=true}.
El código deshabilita el acceso local de desarrollo cuando el entorno
es producción y aplica las verificaciones de autenticación y autorización.
El valor de \texttt{APP\_URL} debe corresponder al dominio publicado,
porque se utiliza al construir los enlaces temporales de entrega.

Las contraseñas, la clave de firma del JWT y las claves secretas de
Supabase permanecen en el backend. Sus valores no se incorporan al
código de Angular ni a esta monografía.

\subsection{Dirección pública y verificación}

\textbf{Dirección pública del sistema:}

\url{https://proyecto-diplomado-angular-frontend.vercel.app}

\textbf{Ruta de salud de la API:}

\url{https://proyecto-diplomado-angular-frontend.vercel.app/api/v1/salud}

La ruta de salud ejecuta una consulta a la base de datos antes de
responder. Durante la verificación documentada en la sección de pruebas,
esta ruta devolvió HTTP 200. También se comprobaron el inicio de sesión,
el rechazo del acceso sin autenticación y la consulta de solicitudes
radiográficas. Las evidencias se conservan en
\path{evidencia/verificacion-despliegue.json}.

Estas comprobaciones corresponden a la versión desplegada en el momento
de la medición. Los cambios de autorización y las rutas específicas
incorporados durante esta revisión se verificaron localmente y requieren
su publicación y comprobación sobre la misma revisión del código.
La entrega completa mediante carga directa a Supabase Storage desde
Vercel no se presenta como validada en esta revisión.
